\section{Part IX --- Rehearsal: questions you will be asked}

\newcommand{\Q}[1]{\par\medskip\noindent\textbf{Q. #1}\par\nobreak\noindent\textbf{A.} }

\Q{What exactly did you change in PhyloGFN?}
Nothing in the mathematics of the sampler --- trajectory balance, uniform backward policy, importance-weighted marginal likelihood --- and one thing in the state space: a second action, \textsc{Reticulate}. Everything else follows from that one addition: the state must know which reticulations are still open (masks), the features must represent two displayed alternatives (two vectors per lineage), the reward must include the reticulation prior, the policy must score one more kind of action, and the batches must handle trajectories of different lengths.

\Q{Why level-1 and not general networks?}
Three reasons. The likelihood of a level-1 network can be computed in one bottom-up pass (Proposition~1) instead of $2^R$ passes; the legality of each action is decided by a constant-time rule on the two lineages involved; and the class is well understood combinatorially, so we can check exhaustively that the generator reaches exactly it (603 for four taxa, 11\,460 for five).

\Q{How do you know the sampler only produces level-1 networks and all of them?}
By construction (Part~IV.2) and by enumeration: \code{test\_class\_counts\_and\_no\_dead\_ends} performs a depth-first search over every reachable state for $N\le5$ and counts the distinct terminal networks by a canonical string; the numbers equal the published counts of Bouvel, Gambette and Mansouri for $R\le2$. Every reachable non-terminal state has a legal action.

\Q{Why is the two-vector likelihood exact?}
Because Felsenstein's recursion is linear in each child vector and the cycles of a level-1 network are vertex-disjoint. At the merge that closes a cycle, the two displayed alternatives are ``keep slot 1'' and ``keep slot 2''; the feature of each carrier under each alternative is available ($A$ or $B$), the alternatives are weighted by $\theta_h$ and $1-\theta_h$ folded into the $A$ vectors, and the sum of the two products is the $\theta$-mixture of the two families of displayed trees --- for each site separately, as the model requires. Deleting a parent edge and suppressing the node is exactly ``multiply by the all-ones vector'' because $P(b)\mathbf1=\mathbf1$ and $P(b_1)P(b_2)=P(b_1+b_2)$. The test compares against explicit enumeration to $10^{-11}$.

\Q{Why can't the two carriers of a reticulation carry different open reticulations?}
A tree node of a binary network has degree three; two edge-disjoint cycles through it would need four incident edges. So in a level-1 binary network cycles are vertex-disjoint, and forbidding $|\mathrm{open}(u)\cup\mathrm{open}(v)|=2$ loses nothing. The enumeration confirms it.

\Q{Why does \textsc{Reticulate} require another free lineage?}
Otherwise the state made only of bare slots, $\{p_1,p_2,q_1,q_2\}$, is reachable and has no legal action --- a dead end, which a GFlowNet cannot have (every state must lead to a terminal state). With the rule, an inductive argument shows every reachable state has a legal action; the exhaustive search confirms zero dead ends.

\Q{What is the backward policy and why does the parent count matter?}
Uniform over parents, as in PhyloGFN. Trajectory balance needs $\log\PB(\tau\mid x)=-\sum_t\log|\mathrm{Pa}(s_{t+1})|$. A network state has two kinds of parents: undo a merge (one per tree-root lineage) and undo a reticulation (one per reticulation whose slots are both bare) --- but the latter only if the forward \textsc{Reticulate} would have been legal from that parent, i.e.\ if the state has a free lineage. Counting a non-existent parent would make $\PB$ not a distribution over the true parents and silently bias the sampler. The round-trip test checks the counts against the backward sampler.

\Q{Why is the terminal object rooted, when PhyloGFN's trees are unrooted?}
Reticulations have a direction (parents above, child below), so a network is inherently rooted. The only consequence is the last-step parent count (1 instead of $2N-3$) and the topology prior normaliser; with $R_{\max}=0$ the evidence targeted is the same as PhyloGFN's because under a reversible model every rooting of an unrooted tree has the same likelihood.

\Q{How is $\theta$ sampled and where does its density enter?}
At the \textsc{Reticulate} step, from a Gaussian mixture over $\mathrm{logit}\,\theta$ (jointly with $\log b_{hx}$). The transition density of the step is the discrete probability of choosing \textsc{Reticulate}$(x)$ times the density of $(b_{hx},\theta)$, converted with the Jacobians $1/b$ and $1/(\theta(1-\theta))$; this is what enters $\log\PF$ in the trajectory-balance loss. The prior $p(\theta)$ is uniform, so it contributes nothing to $\log\R$.

\Q{What prior do you put on the number of reticulations, and why?}
$p(G)\propto e^{-\lambda_RR(G)}$ on the level-1 class with $R\le R_{\max}$, implemented as $-\lambda_RR$ added to the log-reward. Without it, adding parameters ($\theta$, extra branches) can only increase the maximum likelihood, and the sampler would prefer the largest $R$. The posterior over $R$ is read off the histogram of sampled $R$ in the evaluation; the proposal's hyper-prior on $\lambda_R$ is the planned refinement.

\Q{Your negative control fails: on tree-only data the sampler returns $R=2$. Is the method wrong?}
The method is not wrong; the prototype is not converged and the prior is not finished --- and I can show both with numbers. The best $R=2$ network in the buffer beats the re-optimised true tree by 12 nats of \emph{density}, which is what six near-zero extra edges under an exponential prior with $\lambda=10$ ($\log10=2.3$ nats each) buy. The posterior over $R$ integrates that density over the extra dimensions and removes the effect; a GFlowNet that satisfies trajectory balance reproduces the integral, but mine was trained for ten minutes and its Pearson correlation is 0.5--0.7 instead of the paper's 0.99. The fix is on the validation plan: train with the paper's schedule, check against the exact posterior at five taxa, and set the reticulation prior from that check --- the hyper-prior on $\lambda_R$ is scheduled for exactly this reason.

\Q{Your sampler prefers $R=2$ on data simulated with one reticulation. Is it broken?}
No --- it is sampling the posterior under the prior it was given. With $p(G)\propto e^{-\lambda_RR}$ over the whole class and $\lambda_R=1$, the prior mass of $R=2$ exceeds that of $R=1$ by a factor $r(6,2)e^{-2}/(r(6,1)e^{-1})=1.8$ and that of $R=0$ by $15\times$: the class-size growth outweighs the penalty. Switching to \code{UNIFORM\_R} ($p(R)\propto e^{-\lambda_RR}$, uniform within each $R$) removes this, and the negative control (tree-only data) is the test that tells the two apart. This is precisely why the validation plan has a negative control (T8) and why the prior on $R$ needs a hyper-prior or an explicit per-$R$ form rather than a hand-set $\lambda_R$.

\Q{How do you estimate the marginal likelihood for networks?}
Exactly PhyloGFN's importance-weighted bound (eq.~5) with the topology prior $e^{-\lambda_RR}/Z_\lambda$; $Z_\lambda$ uses the exact number of level-1 networks per $R$ from a published closed formula, cross-checked against the paper's tables and our enumeration.

\Q{What does the tree-mode check show?}
With $R_{\max}=0$ the network code is a tree sampler with PhyloGFN's reward; trained for the same small budget on the same alignment, it reaches an MLL within about two log units of the original implementation (both are lower bounds of the same quantity; the difference is training noise at this budget, not a modelling difference).

\Q{Where are the two phases of your proposal in this code?}
Interleaved, as in PhyloGFN: every step attaches its own continuous values (lengths, $\theta$). The terminal object $(G,\bb,\btheta)$ and its reward are identical to the two-phase description; only the intermediate states differ. Both are valid GFlowNet state graphs.

\Q{What is the cost of a step and of the reward now?}
Policy: $\binom{l}{2}+l$ candidate actions, $l\le N+R_{\max}$ --- the same $O(N^2)$ term the proposal's $k$-NN filtering targets. Reward: one Felsenstein pass with at most two vectors per lineage, $O(N\,m\,|\Sigma|^2)$ --- linear in $R$, not $2^R$.

\Q{What would break if a merge joined the two bare slots of the same reticulation?}
Two parallel edges into $h$ and a cycle with no other node: not a binary level-1 network. It is masked; in the backward direction no tree node ever has both children equal to the same reticulation, so the count is consistent.

\Q{Why are there two versions even for a bare slot?}
Uniformity of the update rule: a bare slot is a lineage like any other, with $A=\theta L_h$ (``this side keeps its edge into $h$'') and $B=\mathbf 1$ (``the edge is deleted, nothing hangs here''). The merge rule then never needs a special case for slots.

\Q{How did you make sure you did not break the original?}
The original files are untouched except for four registrations; \code{ENVIRONMENT\_TYPE: ONE\_STEP\_BINARY\_TREE} still runs the original code path, and \file{tools/compare\_tree\_mode.py} runs it side by side with the new one.

\Q{What are the next three things to do?}
(1) Batch the four merge cases to recover PhyloGFN's speed; (2) run DS1--DS8 with $R_{\max}\in\{1,2\}$ and report the posterior over $R$ (tree-only benchmarks should give mode $R=0$: the negative control T8); (3) add the per-locus mixture and the hyper-prior on $\lambda_R$.
\clearpage
